\documentclass[book.tex]{subfiles}
\begin{document}

\chapter{Kohn-Sham Equations as Regularizers}
\label{chap:dft_e3nn}
\noindent\rule{11cm}{0.4pt} \\
\textbf{Prerequisites:} \autoref{chap:molecular_hamiltonian}, \autoref{chap:dft} \\
\textbf{Difficulty Level:} ** \\
\noindent\rule{11cm}{0.4pt}
\vspace{5mm}

%\url{https://journals.aps.org/prl/abstract/10.1103/PhysRevLett.126.036401}

%"""
%Including prior knowledge is important for effective machine learning models in physics, and is usually achieved by explicitly adding loss terms or constraints on model architectures. Prior knowledge embedded in the physics computation itself rarely draws attention. Li will show that solving the Kohn-Sham equations when training neural networks for the exchange-correlation functional provides an implicit regularization that greatly improves generalization. Two separations suffice for learning the entire one-dimensional H2 dissociation curve within chemical accuracy, including the strongly correlated region. The models also generalize to unseen types of molecules and overcome self-interaction errors. Li will also do a code-walkthrough session following the main talk. 
%"""

A basic idea when applying differentiable programming techniques to density functional theory is to use a neural network to approximate the exchange correlation functional. A vanilla approximation though is underspecified and will require large amounts of training data. We can adapt an idea from techniques like physics inspired neural networks though and use the Kohn-Sham equations to constrain the form of the function we seek to learn. This technique is called a Kohn-Sham regularizer \cite{li2021kohn},  \cite{kalita2021generalizability}. The chosen regularization scheme allows for learned models to generalize to new types of molecules past those in a training distribution due to the richer inductive prior.

\section{Kohn-Sham as Regularizer}

The Kohn-Sham equations are given by

\begin{align}
    \left [ -\frac{\nabla^2}{2} + v_s[n](r) \right] \phi_i(r) &= \epsilon_i \phi_i(r)
\end{align}
Here $n$ is the electron density
\begin{align}
    n(r) &= \sum_i |\phi_i(r)|^2
\end{align}
% (\textcolor{red}{TODO: Define a molecular orbital in the Molecular Hamiltonian chapter}).
and $\phi_i$ is a molecular orbital. $v_s$ is the Kohn-Sham potential and is given by
\begin{align}
v_s[n](r) &= v(r) + v_H[n](r) + v_{xc}[n](r)
\end{align}
Here $v$ is the external one-body potential, $v_H$ is the Hartree potential. 
%(\textcolor{red}{TODO: Define elsewhere}).

\begin{align}
    v_{xc}[n](r) &= \frac{\delta E_{xc}}{\delta n(r)}
\end{align}
where $E_{xc}$ is the exchange correlation energy
\begin{align}
    E_{xc}[n] &= \int \epsilon_{xc}[n](r) n(r) dr
\end{align}
where $\epsilon_{xc}[n](r)$ is the exchange correlation energy per electron. This quantity is approximated by a neural network
\begin{align}
    \epsilon_{xc,\theta}[n]
\end{align}
The eigenvalue problem for the Kohn-Sham equation is solved by iterating the self-consistent field equations. %(\textcolor{red}{TODO: Make sure to detail in the preliminaries}).
The loss function is given by the deviation from a reference code
\begin{align}
L(\theta) &= \mathbb{E}_{train}\left [\int  (n_{ks}-n_{DMRG})^2/N_e dx \right ] + \mathbb{E}_{train}\left [\sum_{k=1}^K w_k (E_k - E_{DMRG})^2/N_e\right ]
\end{align}

We use global convolution layers as in the Schnet architecture we studied earlier in the book

\begin{align}
G(n(x), \xi_p) &= \frac{1}{2\xi_p} \int  n(x') \exp(-|x-x'|/\xi_p) dx'
\end{align}

%Unlike the approach we learned about in the previous 
%\section{One Dimensional Example}

%\textcolor{red}{TODO: Add in 1D example from appendix of Li Li paper}
%\textcolor{red}{TODO: A discussion is needed of density matrix renormalization group techniques here.}

%\textcolor{red}{TODO: The Li Li paper enforces a symmetry constraint on its integrals}

%\textcolor{red}{TODO: Add discussion of stop gradient to the differentiable programming part.}

%\textcolor{red}{TODO: I think a major limitation of the Li Li work is the dependence on DMRG in the training loss. This means that the method can't be trained on systems where DMRG is not available.}


\end{document}